\section{Conclusions and Future Developments}\label{sect:ConclusionFutureDevelopments}

\subsection{Project Outcome}\label{subsect:ProjectOutcome}

TrainHub began with a practical problem: gym members and professionals often use separate tools for activities that belong to the same relationship. Workout notes, diet plans, appointments, and messages are scattered across different services. The project brought them together in a single Progressive Web App while preserving an interface suited to each role.

Members can access the gym through a rotating QR badge, follow or create a workout plan, record a live session, consult nutrition information, book appointments, and communicate with their professional. From the other side, professionals can prepare plans, review progress, organize appointments, manage clients, and scan access badges.

What sits underneath is one codebase and one database rather than two products sharing a name: two route trees over the same application shell, twenty-one related tables, every write routed through a checked database procedure, reads that survive the loss of the network, writes that pause and replay when it returns, Web Push, and camera-based facility access. The workout half, the part the user research placed first, was rebuilt during the project around repeating weekly runs, so a plan is something a member trains again rather than something that finishes once.

These actions are connected rather than simply collected in one place. A plan becomes available to the member, a completed workout becomes useful information for the professional, and each update leads back to its relevant context. The two experiences remain distinct, but support the same working relationship: members keep autonomy in their routine, while professionals gain a clearer view of what happens between appointments.

\subsection{Known Limitations}\label{subsect:KnownLimitations}

The delivered version is complete for the scope it set, and the following limits are known rather than discovered late. Stating them is more useful than presenting the project as finished in every direction.

\begin{itemize}[leftmargin=*,nosep]
    \item \textbf{Badge redemption checks the role, not the relationship:} the database confirms that whoever scans a badge is a professional, but not that the member is one of their clients. This matches a front desk, where whoever is on duty admits everyone, and it would have to change if scanning ever meant something more than opening the door;
    \item \textbf{Membership expiry is evaluated, never swept:} no scheduled job rewrites a membership to expired when its end date passes. Every check compares the date instead, so behaviour is correct, but the stored status and the effective one can read differently for the same account;
    \item \textbf{Exercise media are placeholders:} producing or licensing a coherent set of images and videos was outside the scope, and borrowing inconsistent material would have damaged the visual coherence the design work established;
    \item \textbf{Delivery is a tunnel, not a deployment:} the application is served over a fixed HTTPS tunnel from a local production build, which satisfies installability and push but is not a hosted service with its own monitoring and release process.
\end{itemize}

\subsection{Future Developments}\label{subsect:FutureDevelopments}

The next steps follow directly from the limitations above and from features that were deliberately postponed:

\begin{itemize}[leftmargin=*,nosep]
    \item \textbf{Assignment-aware check-in:} verify the professional--member relationship on redemption, or introduce a separate front-desk role for gyms where the two are not the same person;
    \item \textbf{Progress for the member:} the progress views exist only on the professional's side. The same data would support a member-facing history of runs, volumes, and streaks;
    \item \textbf{Plan history:} superseded plans are preserved in the database but no screen reads them. A history view would make the versioning visible to the people it protects;
    \item \textbf{Exercise media:} a licensed, consistent set of images or short videos, with the placeholders replaced rather than supplemented;
    \item \textbf{Gym integration:} connect billing, opening hours, and staff availability, so that membership status stops being simulated by a professional's action;
    \item \textbf{Deployment and auditing:} a hosted release with a repeatable build, followed by a performance and accessibility audit on the deployed origin.
\end{itemize}
